% Phase 2 --- Demand and Operating Profit
% Governed by: paper/spec/MAH_model_rebuild_execution_note_for_Codex.md (Sec. 7)
% Status: Phase 2, replayed onto the correct base review-model@cd5b6e4.
% The derivation itself is unchanged from the original approval; only
% subsec:crosswalk-note below is corrected to reference the real manuscript's
% legacy objects (R_i^event, bar R_i^E, Gamma_i) instead of an earlier,
% mistaken reference to a different PDF that turned out not to be the
% canonical current manuscript. See paper/model_notes/06_open_issues_log.md
% OI-01/OI-02.
% Builds only on objects defined in Phase 1
% (paper/model_notes/01_symbols_and_objects.md). Introduces A, epsilon, p,
% y(p;q), c, pi(q,c), phi, beta, R(q,c). Does not introduce any route,
% manufacturing-technology, or MAH-friction object (RL-19).

\section{Demand and Operating Profit}
\label{sec:demand-profit}

\subsection{Goal}

This module replaces the reduced-form market-return objects used in the
pre-existing manuscript --- $R_i^{event}$ (route-independent
commercialization-event return) and $\bar R_i^E$ (entrusted-route return
kernel), both exogenous primitives in
\texttt{paper/manuscript/mah\_route\_indicator\_friction\_model.tex}
eq.~(eq:GI)--(eq:GE) --- with an operating-profit object derived from an
explicit demand curve and a marginal cost. The marginal cost $c$ is left
as a free argument here: it is not yet a technology-specific object.
Phase 3 instantiates $c$ as $c_I(m,k_i)$ or $c_E(m)$; using an
undetermined $c$ here is a deliberate scope choice matching the note's
own Phase 2 specification, not a violation of the no-forward-reference
discipline (RL-19), because $c$ is a free argument of the function
$R(q,c)$, not a claim about any specific route's cost.

\subsection{Baseline demand}
\label{subsec:demand}

Use a minimal constant-elasticity residual demand for a project of value
$q$:
\begin{equation}
y(p;q) = A q\, p^{-\varepsilon}, \qquad \varepsilon > 1.
\label{eq:demand}
\end{equation}
Here $A>0$ is an exogenous market-size shifter, $q>0$ is the project
value index fixed in Phase 1, and $\varepsilon>1$ is the demand
elasticity. \textbf{MAH must not shift $A$, $q$, or $\varepsilon$}
(RL-08): nothing in this module makes any of the three a function of
$M$.

\subsection{Static pricing problem}
\label{subsec:pricing}

Given a marginal manufacturing cost $c>0$, the seller solves
\begin{equation}
\max_{p \geq c} \ \Pi(p) \equiv (p-c)\,y(p;q) = Aq\left(p^{1-\varepsilon} - c\,p^{-\varepsilon}\right).
\label{eq:pricing-problem}
\end{equation}

\paragraph{FOC.} Differentiating,
\[
\Pi'(p) = Aq\left[(1-\varepsilon)p^{-\varepsilon} + c\varepsilon p^{-\varepsilon-1}\right]
= Aq\,p^{-\varepsilon-1}\big[(1-\varepsilon)p + c\varepsilon\big].
\]
Since $Aq\,p^{-\varepsilon-1}>0$ for $p>0$, the interior FOC is
$(1-\varepsilon)p+c\varepsilon=0$, giving
\begin{equation}
p^*(c) = \frac{\varepsilon}{\varepsilon-1}\,c.
\label{eq:pstar}
\end{equation}

\paragraph{SOC.} Differentiating again,
\[
\Pi''(p) = Aq\,p^{-\varepsilon-2}\Big[\varepsilon(\varepsilon-1)p - c\varepsilon(\varepsilon+1)\Big].
\]
At $p=p^*(c)$, $\varepsilon(\varepsilon-1)p^* = \varepsilon^2 c$, so the bracket equals
$\varepsilon^2c - c\varepsilon(\varepsilon+1) = -\varepsilon c <0$. Hence
\begin{equation}
\Pi''(p^*(c)) = Aq\,\big(p^*(c)\big)^{-\varepsilon-2}(-\varepsilon c) < 0,
\label{eq:soc}
\end{equation}
since $A,q,\varepsilon,c,p^*(c)>0$. The interior critical point is a strict
maximum.

\subsection{Operating profit}
\label{subsec:profit}

Substituting \eqref{eq:pstar} into \eqref{eq:pricing-problem}:
\[
\pi(q,c) \equiv \Pi(p^*(c))
= Aq\,\big(p^*(c)\big)^{-\varepsilon}\big(p^*(c) - c\big)
= Aq\,\big(p^*(c)\big)^{-\varepsilon}\cdot \frac{c}{\varepsilon-1},
\]
and using $\big(p^*(c)\big)^{-\varepsilon} = \left(\frac{\varepsilon-1}{\varepsilon}\right)^{\varepsilon} c^{-\varepsilon}$,
\begin{equation}
\boxed{\ \pi(q,c) = Aq\,\frac{(\varepsilon-1)^{\varepsilon-1}}{\varepsilon^{\varepsilon}}\,c^{1-\varepsilon}. \ }
\label{eq:profit}
\end{equation}

\paragraph{Comparative statics.}
\[
\pi_q = \frac{\partial \pi}{\partial q} = A\,\frac{(\varepsilon-1)^{\varepsilon-1}}{\varepsilon^{\varepsilon}}\,c^{1-\varepsilon} = \frac{\pi(q,c)}{q} > 0,
\]
\[
\pi_c = \frac{\partial \pi}{\partial c} = (1-\varepsilon)\,Aq\,\frac{(\varepsilon-1)^{\varepsilon-1}}{\varepsilon^{\varepsilon}}\,c^{-\varepsilon} = (1-\varepsilon)\,\frac{\pi(q,c)}{c} < 0,
\]
where the sign of $\pi_c$ follows because $\varepsilon>1\Rightarrow(1-\varepsilon)<0$ and $\pi(q,c)/c>0$. Both signs match the note's requirement $\pi_q>0,\ \pi_c<0$.

\subsection{Present value}
\label{subsec:pv}

Suppose a commercialized drug remains commercially active into the next
period with probability $\varphi\in[0,1)$, and future profit is
discounted at $\beta\in(0,1)$. The flow of profits is
$\pi,\ \beta\varphi\pi,\ (\beta\varphi)^2\pi,\dots$, a geometric series
with ratio $\beta\varphi\in[0,1)$, so it converges to
\begin{equation}
\boxed{\ R(q,c) = \frac{\pi(q,c)}{1-\beta\varphi}. \ }
\label{eq:pv}
\end{equation}
Because $\beta\varphi<1$ is guaranteed by $\beta<1$ and $\varphi<1$, the
denominator is strictly positive and $R(q,c)$ is well-defined. Since
$1-\beta\varphi>0$ is a positive scalar independent of $q,c$,
\[
R_q = \frac{\pi_q}{1-\beta\varphi} > 0, \qquad
R_c = \frac{\pi_c}{1-\beta\varphi} < 0,
\]
i.e.\ the present-value return inherits the sign pattern of operating
profit exactly.

\subsection{Dimensional consistency}
\label{subsec:dimensions}

All objects in this module are flow-profit quantities per period, except
$\beta,\varphi,\varepsilon$ which are dimensionless scalars and $A$ which
absorbs whatever units are needed to make \eqref{eq:demand} a quantity.
$\pi(q,c)$ has revenue-flow units (price $\times$ quantity); dividing by
the dimensionless $(1-\beta\varphi)$ leaves $R(q,c)$ in the same
revenue-flow units, capitalized into a present value. No route-specific
fixed cost (e.g.\ a future $F_I,F_E,\mu_E,\tau_E$) is introduced or
subtracted in this module, so no cost can be double-counted here; that
check is repeated when route values combine $R(q,c)$ with route-specific
costs in Phase 4.

\subsection{Economic interpretation}
\label{subsec:econ-interp}

Operating profit $\pi(q,c)$ is the constant-markup surplus a seller
retains after pricing optimally against a project-specific demand curve
scaled by $q$ and producing at constant marginal cost $c$: a
higher-value project ($q$ larger) scales profit proportionally at any
given cost, while a higher marginal cost $c$ lowers profit both directly,
through a thinner per-unit markup, and indirectly, through the
higher optimal price it induces, which shrinks quantity sold. $R(q,c)$
capitalizes this per-period profit into a present value under a
survival-hazard structure: $\varphi$ is the chance the product still
generates profit next period, so the effective discount factor
$\beta\varphi$ compounds patience with commercial durability. Because
$A$, $q$, and $\varepsilon$ are demand-side and clinical-value-side
primitives that this module deliberately keeps free of $M$, any MAH
effect on a developer's payoff must enter later, and only through the
cost argument $c$ once Phase 3 gives $c$ technology-specific content —
this is what makes the eventual MAH channel a manufacturing/organizational
one rather than a demand or clinical-value one.

\subsection{Crosswalk to the pre-existing manuscript (completed formally in Phase 12)}
\label{subsec:crosswalk-note}

$R(q,c)$, together with the route-specific costs added in Phase 3--4, is
intended to fully replace the pre-existing manuscript's two exogenous
return primitives (\texttt{paper/manuscript/mah\_route\_indicator\_friction\_model.tex}
\S4.2, just before eq.~(eq:GI)):
\begin{itemize}
\item $R_i^{event}$ --- ``the route-independent commercialization-event
  return from an internally realized original drug'' --- replaced by
  $R(q,c_I(m,k_i))$ once Phase 3 instantiates the internal-route cost.
\item $\bar R_i^E$ --- ``the corresponding entrusted-route return
  kernel'' --- replaced by $R(q,c_E(m))$ once Phase 3 instantiates the
  entrusted-route cost.
\end{itemize}
Both legacy objects are firm-level exogenous constants with no demand or
cost microfoundation; $R(q,c)$ is derived and project-level ($q$-indexed),
which is the intended generalization. The manuscript's inclusive value
$\Gamma_i(M,\eta,p_m)$ (its eq.~(eq:Gamma)) is a separate object --- an
expectation over a \emph{logit} route choice --- and is not replaced
here; it is replaced by $\Omega_i(M,p_m)$ in Phase 5, once Phase 4 route
values exist. No reduced-form return object may survive past this module
without one of these two pointers; the full dependency accounting
(equations, propositions, calibration claims that depended on
$R_i^{event}$ or $\bar R_i^E$) is completed formally in the Phase 12
crosswalk table, per the note's own scoping of that phase.

\subsection{What this module deliberately excludes}

No route value, no manufacturing-cost technology $c_I,c_E,F_I,F_E$, no
MAH friction $\tau_E(M)$, and no advancement/cost object $x_i,C_X$
appears in this module. $M$ does not appear at all in this module: the
demand and profit chain is entirely policy-free, as required (RL-08).
